\subsection{Fundamentación}

El desarrollo de este software surge como respuesta a las problemáticas detectadas en la gestión cotidiana de los edificios, representando a su vez una oportunidad invaluable para aplicar y consolidar los principios de la ingeniería de software en un entorno real. La dependencia de procesos manuales genera demoras en la recaudación y falta de transparencia. Se vuelve imperativo proponer soluciones tecnológicas que reemplacen prácticas obsoletas, como los cuadernos físicos de guardia, que son vulnerables a pérdidas y alteraciones. CoConsorcio se justifica técnica y académicamente al plantear una herramienta acotada pero robusta, construida bajo estándares modernos de desarrollo web, que moderniza el flujo de trabajo del administrador y empodera al vecino. Al tratarse de un proyecto gestado en el ámbito universitario, su diseño prioriza el código limpio, la escalabilidad y la correcta aplicación de metodologías ágiles por sobre la saturación de funcionalidades.
